% Adapted from an OpenAI Math manuscript. See PROVENANCE.md.
\section{Interior lattice points and interpolation}
\label{sec:saturation}

Reshaping gives the right number of exponents but leaves their exact positions unknown. We need an entire simplex of exponents for interpolation. The next lemma obtains interior lattice points from additivity and asymptotically full density. It is a direct argument in the setting of the more general semigroup approximation theorem of \textcite[Theorem~1.6]{KK12}.

\begin{lemma}
\label{lem:interior-saturation}
Let $n\ge1$ be an integer, let $P\subset\mathbb{R}^n$ be a compact convex set with nonempty interior, and let $F_k\subset kP\cap\mathbb{Z}^n$ for every integer $k\ge1$. Suppose that
\[
 F_k+F_l\subset F_{k+l}\quad(k,l\ge1),
 \qquad
 \#F_k=\vol(P)k^n+o(k^n).
\]
For every compact set $K\subset\operatorname{int}P$, there is an integer $k_K$ such that
\[
 kK\cap\mathbb{Z}^n\subset F_k\qquad(k\ge k_K).
\]
\end{lemma}

\begin{proof}
The number of lattice points of $kP$ missing from $F_k$ is $o(k^n)$, by the lattice-point asymptotic for convex bodies. Put $k_1=\lfloor k/2\rfloor$ and $k_2=k-k_1$. If $y\in kK\cap\mathbb Z^n$, the region $k_1P\cap(y-k_2P)$ contains the ball of radius $\delta k_1$ about $k_1y/k$, with $\delta>0$ chosen so that $K+\delta\overline B(0,1)\subset P$. This ball contains at least $ck^n$ lattice points for large $k$, uniformly in its center, as is seen by inscribing a cube. Only $o(k^n)$ of these points are missing from $F_{k_1}$ or have their translates $y-y_1$ missing from $F_{k_2}$. Thus $y\in F_{k_1}+F_{k_2}\subset F_k$.
\end{proof}

In the stretched simplex below, each transverse monomial leaves exactly enough degrees in the last coordinate to interpolate its coefficient independently.

\begin{lemma}
\label{lem:simplex-interpolation}
Let $n\ge2$, $r\ge1$, and $m\ge0$ be integers, and let $c_1,\ldots,c_r\in\mathbb{C}^*$ be distinct. Put
\[
 E_{r,m}=
 \left\{(\beta,b)\in\mathbb{Z}_{\ge0}^{n-1}\times\mathbb{Z}_{\ge0}:
       |\beta|+\frac br<m+1\right\}.
\]
Then evaluation on jets through degree $m$ at
\[
 q_i=(1,\ldots,1,c_i)\in\mathbb{T}^n,\qquad 1\le i\le r,
\]
is an isomorphism from $\mathcal{M}(E_{r,m})$ to the direct sum of the jet spaces.
\end{lemma}

\begin{proof}
Write the torus coordinates as $(x',Y)=(x_1,\ldots,x_{n-1},Y)$ and put $X=x'-(1,\ldots,1)$. Translation from $x'$ to $X$ preserves the space of polynomials with the indicated support. In fact, the coefficient space of $Y^b$ consists of the polynomials of total degree less than $m+1-b/r$ in the first $n-1$ variables. Consequently every element of the space has a unique expression
\[
 \sum_{|\alpha|\le m}X^\alpha f_\alpha(Y),
 \qquad
 \deg f_\alpha<r(m+1-|\alpha|).
\]
At $q_i$ the local parameters are $X$ and $Y-c_i$. For each $\alpha$ the jet records the remainder of $f_\alpha$ modulo $(Y-c_i)^{m+1-|\alpha|}$. As $i$ varies, these powers are pairwise comaximal. The Chinese remainder theorem gives a bijection between the simultaneous remainders and the polynomials of degree less than the degree $r(m+1-|\alpha|)$ of their product. Taking this bijection for each $\alpha$ proves the assertion.
\end{proof}

The two lemmas now convert asymptotically full additive supports into simultaneous jet surjectivity. Translating a slightly smaller simplex into the interior avoids any assertion about exponents on the boundary.

\begin{proposition}
\label{prop:asymptotic-jets}
Let $n\ge2$ and $r\ge1$ be integers, let $w>0$, and set
\[
 P=\Delta(w,\ldots,w,rw).
\]
Suppose that finite sets $F_k\subset kP\cap\mathbb{Z}^n$, indexed by all integers $k\ge1$, satisfy
\[
 F_k+F_l\subset F_{k+l},
 \qquad
 \#F_k=\vol(P)k^n+o(k^n).
\]
For every $0<\theta<w$ and every choice of distinct $c_1,\ldots,c_r\in\mathbb{C}^*$, the space $\mathcal{M}(F_k)$ surjects onto jets through degree $\lfloor k\theta\rfloor$ at $(1,\ldots,1,c_i)$ for all sufficiently large $k$.
\end{proposition}

\begin{proof}
Choose $\theta/w<\lambda<1$. There is a vector $a$ with strictly positive coordinates such that $a+\lambda P\subset\operatorname{int}P$: it suffices to require
\[
 \frac{a_1+\cdots+a_{n-1}}w+\frac{a_n}{rw}<1-\lambda.
\]
Choose $a_k\in k^{-1}\mathbb{Z}^n$ with $a_k\to a$. For all sufficiently large $k$, the translated simplices $a_k+\lambda P$ lie in a single compact subset $K$ of $\operatorname{int}P$. \Cref{lem:interior-saturation} therefore gives
\[
 ka_k+(k\lambda P\cap\mathbb{Z}^n)\subset F_k.
\]
Consequently $\mathcal{M}(F_k)$ contains the monomial space for $k\lambda P\cap\mathbb{Z}^n$, multiplied by the monomial with exponent $ka_k$. Multiplication by this monomial is an isomorphism on each jet algebra at a torus point, since the monomial is a unit there.

Put $m=\lfloor k\theta\rfloor$. For all sufficiently large $k$, $m+1\le k\lambda w$, and hence
\[
 E_{r,m}\subset k\lambda P\cap\mathbb{Z}^n.
\]
\cref{lem:simplex-interpolation} supplies jet surjectivity for this subspace. The monomial translation preserves surjectivity and proves the claim.
\end{proof}
